\section{The \texttt{Ring\_p} McXtrace Component}
Date: Feb. 2017
Version: 1.1
Release: McXtrace 1.2
Origin: DTU Physics, DTU Space

Stack of conical shells as part of a Wolter optic. Parabolic version.

\subsection*{Identification}
\begin{itemize}
  \item \textbf{Author:} Erik B Knudsen and Desiree D. M. Ferreira
  \item \textbf{Origin:} DTU Physics, DTU Space
  \item \textbf{Date:} Feb. 2016
\end{itemize}

\subsection*{Description}
A stack of conical shells is simulated. Parabolic version. To intersect the Wolter I plates we take advatage of the azimuthal symmetry and only consider the radial component of the photon's wavevector.

Example: Ring\_p( pore\_th=0, ring\_nr=3, Z0=FL, yheight=0.605e-3, mirror\_reflec="mirror\_coating\_unity.txt", R\_d=0, size\_file="ref\_design\_breaks.txt")

\subsection*{Input parameters}
Parameters in \textbf{boldface} are required; the others are optional.

\begin{longtable}{p{0.22\textwidth}p{0.12\textwidth}p{0.46\textwidth}p{0.14\textwidth}}
\toprule
\textbf{Name} & \textbf{Unit} & \textbf{Description} & \textbf{Default} \\
\midrule
\endhead
\textbf{ring\_nr} &  &  &  \\
\textbf{pore\_th} &  &  &  \\
\textbf{Z0} & m & Distance between optics centre plane and focal spot (essentially focal length). &  \\
\textbf{yheight} & m & Height of the pore. &  \\
chamferwidth & m & Width of side walls. & 0 \\
gap & m & Gap between the plate and the intersection plane with the hyperbolic section. & 0 \\
zdepth &  &  & 0 \\
mirror\_reflec &   & Data file containing reflectivities of the reflector surface (TOP). & "" \\
bottom\_reflec &   & Data file containing reflectivities of the bottom surface (BOTTOM). & "" \\
size\_file &  &  & "" \\
R\_d &   & Default reflectivity value to use if no reflectivity file is given. Useful f.i. is one surface is reflecting and the others absorbing. & 1 \\
waviness & rad & Waviness of the pore reflecting surface. The slope error is assumed to be uniformly distributed in the interval [-waviness:waviness]. & 0 \\
longw &   & If non-zero, waviness is 1D and along the pore axis. & 0 \\
\bottomrule
\end{longtable}

\subsection*{Links}
\begin{itemize}
  \item Component source code found in file \texttt{Ring\_p.comp}.
\end{itemize}
\IfFileExists{astrox/Ring_p_static.tex}{\input{astrox/Ring_p_static.tex}}{}